\chapter{Arquitetura da Solu\c{c}\~ao e Integra\c{c}\~ao dos Dados}
\label{cap:arquitetura}

Este cap\'itulo apresenta a arquitetura final da solu\c{c}\~ao e a forma como os dados escolares foram integrados para viabilizar a previs\~ao da pr\'oxima nota e a identifica\c{c}\~ao de risco pedag\'ogico. A organiza\c{c}\~ao aqui descrita n\~ao surgiu pronta desde o in\'icio; ela foi sendo consolidada \`a medida que o projeto deixou de ser apenas um conjunto de experimentos e passou a exigir uma estrutura reproduz\'ivel, compreens\'ivel e segura para publica\c{c}\~ao acad\^emica.

\section{Vis\~ao geral da arquitetura}
\label{sec:visao-geral}
\label{sec:visao-geral-solucao}

O fluxo consolidado pode ser resumido em oito etapas:
\begin{enumerate}
    \item manuten\c{c}\~ao e anonimiza\c{c}\~ao da base original;
    \item extra\c{c}\~ao das tabelas relevantes para arquivos \textit{Comma-Separated Values} (CSV);
    \item leitura estruturada dessas fontes pela pipeline;
    \item constru\c{c}\~ao do dataset longitudinal;
    \item separa\c{c}\~ao temporal entre treino, valida\c{c}\~ao e teste;
    \item compara\c{c}\~ao autom\'atica entre modelos e baselines;
    \item gera\c{c}\~ao de predi\c{c}\~oes e an\'alises de erro;
    \item produ\c{c}\~ao de relat\'orios operacionais para a escola.
\end{enumerate}

Em termos de organiza\c{c}\~ao do c\'odigo, o n\'ucleo da solu\c{c}\~ao ficou distribu\'ido entre m\'odulos de configura\c{c}\~ao, leitura das entradas, constru\c{c}\~ao do dataset, modelagem, gera\c{c}\~ao de artefatos e relat\'orios. A Figura~\ref{fig:arquitetura-solucao} sintetiza essa organiza\c{c}\~ao em camadas maiores, priorizando a leitura do caminho dos dados e das responsabilidades principais da solu\c{c}\~ao.

\begin{figure}[H]
\caption{Vis\~ao arquitetural da solu\c{c}\~ao proposta.}
\label{fig:arquitetura-solucao}
\centering
\includeprojectdiagram[width=0.95\textwidth]{fig/diagrama-export/arquitetural.pdf}
\fontefig{Elaborada pelo autor}
\end{figure}

\section{Modos de execu\c{c}\~ao}
\label{sec:modos}

A aplica\c{c}\~ao foi organizada para funcionar em dois cen\'arios complementares. No cen\'ario completo, a rotina parte da prepara\c{c}\~ao local do banco restaurado, segue para a extra\c{c}\~ao dos CSVs can\^onicos e termina na execu\c{c}\~ao da pipeline com gera\c{c}\~ao de relat\'orios. No cen\'ario p\'ublico reproduz\'ivel, o ponto de partida s\~ao os CSVs j\'a tratados e anonimizados, a partir dos quais os modos anal\'iticos e a recomposi\c{c}\~ao da camada final podem ser executados sem acesso ao banco institucional.

Esse desenho permite separar o que depende de infraestrutura local controlada do que pode ser demonstrado, documentado e reproduzido academicamente. A Figura~\ref{fig:fluxo-pipeline} representa o fluxo operacional principal, desde a prepara\c{c}\~ao e extra\c{c}\~ao at\'e a execu\c{c}\~ao dos modos de modelagem e a gera\c{c}\~ao dos relat\'orios finais.

\begin{figure}[H]
\caption{Fluxo operacional principal da solu\c{c}\~ao.}
\label{fig:fluxo-pipeline}
\centering
\includeprojectdiagram[width=0.95\textwidth]{fig/diagrama-export/fluxo.pdf}
\fontefig{Elaborada pelo autor}
\end{figure}

\section{Configura\c{c}\~ao experimental e reprodutibilidade}
\label{sec:config-experimental}

\subsection{Camadas de entrada}

A reprodu\c{c}\~ao da solu\c{c}\~ao foi separada em duas camadas complementares:
\begin{itemize}
    \item camada privada local: prepara o banco restaurado, aplica anonimiza\c{c}\~ao e executa as consultas SQL reais fora do Git;
    \item camada p\'ublica reproduz\'ivel: consome os CSVs can\^onicos em \path{artifacts/database/csv/} e executa a pipeline acad\^emica publicada no reposit\'orio.
\end{itemize}

Essa separa\c{c}\~ao foi adotada para que a parte sens\'ivel da extra\c{c}\~ao permanecesse protegida, enquanto a parte acad\^emica e demonstrativa da solu\c{c}\~ao continuasse pass\'ivel de reprodu\c{c}\~ao por terceiros.

\subsection{Arquivos de entrada can\^onicos}

Os arquivos de entrada da arquitetura final foram definidos da seguinte forma:
\begin{itemize}
    \item \texttt{aluno.csv}: identifica o aluno de forma pseudonimizada e preserva o contexto cadastral e escolar b\'asico da extra\c{c}\~ao;
    \item \texttt{media\_nota\_aluno.csv}: base principal de notas por aluno, disciplina, per\'iodo e etapa; \'e o eixo central da modelagem;
    \item \texttt{faltas\_aluno.csv}: registros de faltas usados para formar faltas por etapa e faltas acumuladas;
    \item \texttt{pagamento\_aluno.csv}: registros financeiros resumidos para produzir sinais administrativos por aluno e ano letivo;
    \item \texttt{responsaveis\_aluno.csv}: arquivo can\^onico complementar do projeto, preservado como parte do contrato p\'ublico da extra\c{c}\~ao;
    \item \texttt{professor\_disciplina.csv}: contexto docente por disciplina e turma, incorporado quando esse v\'inculo est\'a dispon\'ivel.
\end{itemize}

Na pipeline central publicada, os arquivos consumidos diretamente pela modelagem e pelos relat\'orios s\~ao, principalmente, \path{media_nota_aluno.csv}, \path{faltas_aluno.csv}, \path{pagamento_aluno.csv} e \path{professor_disciplina.csv}. Os demais arquivos preservam o contrato p\'ublico da extra\c{c}\~ao e a rastreabilidade da arquitetura completa, mesmo quando n\~ao entram diretamente em todas as etapas da modelagem.

\subsection{Execu\c{c}\~ao reproduz\'ivel da pipeline}

A execu\c{c}\~ao experimental do projeto pode ser compreendida em dois cen\'arios complementares.

No cen\'ario completo, dispon\'ivel apenas em ambiente local com acesso controlado ao banco restaurado, a rotina segue tr\^es movimentos:
\begin{enumerate}
    \item preparar localmente o banco restaurado por meio da rotina \texttt{prepare-db};
    \item extrair os arquivos can\^onicos por meio da rotina \texttt{extract};
    \item executar a pipeline completa com a rotina \texttt{workflow}, respons\'avel por acionar os dois modos principais e consolidar a camada final de relat\'orios.
\end{enumerate}

No cen\'ario p\'ublico reproduz\'ivel, basta disponibilizar os CSVs can\^onicos em \path{artifacts/database/csv/}. A partir deles, a execu\c{c}\~ao pode ocorrer por modo isolado ou de ponta a ponta:
\begin{itemize}
    \item \texttt{pipeline --mode previsao\_nota}: executa apenas o modo voltado \`a regress\~ao da pr\'oxima nota;
    \item \texttt{pipeline --mode alerta\_risco}: executa apenas o modo voltado \`a classifica\c{c}\~ao de risco;
    \item \texttt{reports}: recomp\~oe a camada final de relat\'orios a partir dos artefatos j\'a produzidos;
    \item \texttt{workflow}: reconstitui o fluxo completo e produz a sa\'ida operacional final.
\end{itemize}

Essa organiza\c{c}\~ao foi importante porque o experimento final n\~ao depende de acesso direto ao banco para ser reproduzido academicamente; ele depende do contrato dos CSVs e de uma interface de execu\c{c}\~ao padronizada.

\subsection{Modos e artefatos anal\'iticos}

Cada modo da pipeline grava seus resultados em um subdiret\'orio de \path{artifacts/pipeline/}. Os diret\'orios centrais s\~ao:
\begin{itemize}
    \item \path{artifacts/pipeline/previsao_nota/};
    \item \path{artifacts/pipeline/alerta_risco/}.
\end{itemize}

Em ambos os casos, os principais artefatos gerados s\~ao:
\begin{itemize}
    \item \texttt{student\_prediction\_dataset.csv}: base final preparada para a rodada de modelagem;
    \item \texttt{student\_prediction\_predictions.csv}: predi\c{c}\~oes e sinais de teste usados na avalia\c{c}\~ao;
    \item \texttt{student\_prediction\_report.txt}: resumo t\'ecnico com m\'etricas e candidatos selecionados;
    \item \texttt{student\_prediction\_model.pkl}: modelo final persistido quando a solu\c{c}\~ao vencedora depende de objeto supervisionado treinado;
    \item \texttt{error\_analysis\_by\_subject.csv}: consolida\c{c}\~ao de erro por disciplina;
    \item \texttt{error\_analysis\_by\_series.csv}: consolida\c{c}\~ao de erro por s\'erie;
    \item \texttt{error\_analysis\_by\_score\_band.csv}: consolida\c{c}\~ao de erro por faixa de nota;
    \item \texttt{error\_analysis\_summary.txt}: resumo textual das an\'alises de erro.
\end{itemize}

Esses artefatos sustentam a leitura t\'ecnica do experimento, pois registram o tamanho da base resultante, os candidatos comparados, o desempenho temporal, as predi\c{c}\~oes no teste e os pontos de maior dificuldade do sistema.

\subsection{Artefatos operacionais finais}

Ap\'os a execu\c{c}\~ao dos dois modos, a camada final de disponibiliza\c{c}\~ao grava sa\'idas em \path{artifacts/reports/}. Entre os arquivos mais relevantes est\~ao:
\begin{itemize}
    \item \texttt{relatorio\_escolar\_integrado.csv}: base consolidada que cruza previs\~ao de nota, alerta de risco e sinais operacionais;
    \item \texttt{relatorio\_professor.csv} e \texttt{relatorio\_professor.txt}: vis\~ao por aluno e disciplina para apoio pedag\'ogico direto;
    \item \texttt{relatorio\_coordenacao.csv} e \texttt{relatorio\_coordenacao.txt}: consolidados por s\'erie, disciplina e prioridade institucional;
    \item \texttt{relatorio\_secretaria.csv} e \texttt{relatorio\_secretaria.txt}: sinais administrativos e prioridades de contato;
    \item \texttt{relatorio\_professores\_por\_turma.csv}: consolida\c{c}\~ao de risco por professor e turma;
    \item \texttt{relatorio\_professores\_geral.csv}: consolida\c{c}\~ao geral de risco por professor;
    \item \texttt{ranking\_executivo\_coordenacao} em formatos \texttt{.csv} e \texttt{.txt}: prioriza\c{c}\~ao executiva dos alunos mais urgentes para interven\c{c}\~ao.
\end{itemize}

Esses relat\'orios representam a sa\'ida final do experimento no contexto do TCC, pois convertem resultados de modelagem em materiais de consulta orientados aos diferentes perfis institucionais da escola.

\subsection{Condi\c{c}\~ao m\'inima de reprodutibilidade}

Assim, a condi\c{c}\~ao m\'inima para reproduzir publicamente a configura\c{c}\~ao experimental n\~ao \'e o acesso ao banco institucional, mas sim a disponibilidade dos CSVs can\^onicos e a observ\^ancia do fluxo operacional da aplica\c{c}\~ao.

\section{Integra\c{c}\~ao de notas, faltas, pagamentos e professores}
\label{sec:fontes-dados-integradas}
\label{sec:integracao}

O trabalho foi estruturado para unir quatro blocos principais de informa\c{c}\~ao:
\begin{itemize}
    \item desempenho acad\^emico;
    \item frequ\^encia;
    \item sinais administrativos e financeiros;
    \item contexto escolar, incluindo disciplina, s\'erie, turma e professor.
\end{itemize}

Do ponto de vista de engenharia, a entrada de dados foi organizada em duas camadas. A primeira camada \'e privada e local, respons\'avel pela prepara\c{c}\~ao do banco restaurado e pela execu\c{c}\~ao das consultas SQL reais. A segunda camada \'e p\'ublica e reproduz\'ivel, baseada nos arquivos CSV can\^onicos consumidos pela pipeline em \path{artifacts/database/csv/}. Isso permite explicar e reproduzir a arquitetura sem divulgar a estrutura interna do banco institucional.

Conforme sintetizado na Se\c{c}\~ao~\ref{sec:config-experimental}, essa separa\c{c}\~ao define o contrato de reprodu\c{c}\~ao da solu\c{c}\~ao. Nesta se\c{c}\~ao, o foco passa a ser a forma como esses blocos de dados foram integrados na modelagem. As notas representam o eixo central da modelagem. A partir delas, foram constru\'idos hist\'oricos deslocados, m\'edias m\'oveis, tend\^encias e estat\'isticas de refer\^encia. As faltas foram agregadas por etapa e acumuladas ao longo do tempo. Os pagamentos foram resumidos por ano letivo, incluindo quantidades, pend\^encias e m\'edias de atraso. O professor, quando dispon\'ivel no banco atualizado, foi incorporado como atributo contextual, sem remover linhas quando o v\'inculo docente estivesse ausente.

\section{Engenharia de atributos}
\label{sec:atributos-temporais-construidos}
\label{sec:engenharia}

Os atributos criados para a modelagem incluem:
\begin{itemize}
    \item nota anterior 1, 2 e 3;
    \item m\'edia das \'ultimas notas;
    \item desvio hist\'orico da disciplina;
    \item tend\^encia da \'ultima nota;
    \item m\'edia hist\'orica do aluno no ano;
    \item m\'edia do aluno no ano anterior;
    \item m\'edia da disciplina no ano anterior;
    \item m\'edia da corte por etapa;
    \item faltas da etapa;
    \item faltas acumuladas;
    \item movimentos financeiros registrados e pendentes no ano;
    \item dias at\'e pagamento;
    \item indicador de disponibilidade do professor.
\end{itemize}

Os atributos diretamente derivados da trajet\'oria de notas e faltas foram constru\'idos por deslocamentos, acumulados ou agrega\c{c}\~oes at\'e a etapa observada, preservando a causalidade temporal do problema. Os sinais financeiros, por sua vez, foram usados como contexto administrativo agregado por aluno e ano letivo. Por isso, devem ser interpretados como vari\'aveis institucionais de apoio, e n\~ao como evid\^encia causal isolada do desempenho.

\section{Constru\c{c}\~ao do alvo e preven\c{c}\~ao de vazamento}
\label{sec:definicao-alvos-dataset}
\label{sec:alvo}
\label{sec:regras-evitar-vazamento}
\label{sec:prevencao-vazamento}

Foram definidos dois alvos principais:
\begin{enumerate}
    \item \texttt{target\_nota\_proxima}, correspondente ao valor num\'erico da pr\'oxima nota observada;
    \item \texttt{target\_risco\_proxima}, correspondente ao indicador bin\'ario de pr\'oxima nota inferior a 6,0.
\end{enumerate}

Para evitar vazamento de informa\c{c}\~ao, os atributos hist\'oricos de desempenho e frequ\^encia foram calculados por deslocamento, m\'edias acumuladas ou agrega\c{c}\~oes limitadas ao passado. Isso foi essencial para garantir que o desempenho medido refletisse um cen\'ario real de uso, e n\~ao apenas um ajuste artificial a dados futuros. A preven\c{c}\~ao de vazamento, neste caso, n\~ao foi um detalhe t\'ecnico menor, mas uma condi\c{c}\~ao para que os resultados tivessem valor pr\'atico.

\section{Camada de relat\'orios}
\label{sec:saidas-geradas-pipeline}
\label{sec:camada-relatorios}

Al\'em da modelagem, a arquitetura incorpora uma camada pr\'opria de relat\'orios. Ela organiza sa\'idas para os principais perfis escolares e consolidados complementares:
\begin{itemize}
    \item relat\'orio do professor, com foco por aluno e disciplina;
    \item relat\'orio da coordena\c{c}\~ao, com agregados por s\'erie e disciplina;
    \item relat\'orio da secretaria, com \^enfase administrativa;
    \item ranking executivo de alunos priorit\'arios;
    \item relat\'orio escolar integrado, usado como base operacional unificada;
    \item consolidados de risco por professor, com e sem recorte por turma.
\end{itemize}

Essa camada foi decisiva para aproximar a solu\c{c}\~ao do problema real da escola, convertendo resultados t\'ecnicos em informa\c{c}\~ao acion\'avel. O invent\'ario operacional desses artefatos foi consolidado na Se\c{c}\~ao~\ref{sec:config-experimental}. Essa escolha tamb\'em ajuda a marcar a diferen\c{c}a entre o que o modelo calcula e o que a escola efetivamente consome.

\section{Estrutura final do reposit\'orio e do fluxo operacional}
\label{sec:organizacao-pacote-school-predictor}
\label{sec:estrutura-final}

Na forma final do projeto, a base ativa do software ficou organizada nos seguintes blocos:
\begin{itemize}
    \item \texttt{school\_predictor/database}: acesso ao banco, wrappers p\'ublicos de prepara\c{c}\~ao e extra\c{c}\~ao, e contrato das consultas;
    \item \texttt{school\_predictor/database/private\_runtime.py}: camada local privada para prepara\c{c}\~ao f\'isica do banco e extra\c{c}\~ao com tratamento sens\'ivel;
    \item \texttt{school\_predictor/database/private\_sql/}: consultas SQL reais mantidas apenas no ambiente local;
    \item \texttt{school\_predictor/pipeline}: constru\c{c}\~ao do dataset, modelagem, avalia\c{c}\~ao e relat\'orios;
    \item \texttt{school\_predictor/app}: dashboard em Streamlit;
    \item \texttt{artifacts/database/csv}: interface can\^onica de entrada da pipeline;
    \item \texttt{artifacts/pipeline} e \texttt{artifacts/reports}: sa\'idas anal\'iticas e operacionais.
\end{itemize}

Com isso, a arquitetura final passou a separar claramente o que precisa ficar privado por envolver estrutura institucional do banco e o que pode ser publicado para fins acad\^emicos, reprodutibilidade e apresenta\c{c}\~ao do TCC. Esse arranjo resume bem a dire\c{c}\~ao adotada pelo projeto: preservar o contexto real da escola sem comprometer a possibilidade de explicar, auditar e continuar a solu\c{c}\~ao.
